\section*{O teste decisivo: a razão de Hubble local/fundo (H0DN $\times$ BG1), a emenda e a abertura}
\addcontentsline{toc}{section}{O teste decisivo: a razão de Hubble local/fundo (H0DN $\times$ BG1), a emenda e a abertura}
\textbf{A lei e a pergunta.} Nulo: $H_{\rm local}/H_{\rm fundo}=1$ (o $\Lambda$CDM). Alternativa: $K=E(z_\star)^{2\beta/3}$, com $\beta=\alpha\sqrt e$ em runtime e zero parâmetro livre ($K=1{,}083807$, $\log_{10}K=0{,}034952$, recomputados neste rito). A pergunta: o número local da rede H0DN, lido contra o fundo BG1, cabe na razão 1 ou na razão $K$?
\textbf{O que a v410 tinha pré-registrado e por que não abriu.} A v410 pré-registrou a função de veredito e não leu dado observado; a ordem do operador era não abrir sem poder discriminante, e o vetor publicado não cabia no suporte da covariância publicada. Com os mesmos números lidos agora, a função congelada da v410 (a do \texttt{um.py}) devolve o ramo (1): \texttt{INCONCLUSIVE\_\allowbreak{}SYSTEMATICS\_\allowbreak{}\_\allowbreak{}RETRODICTION\_\allowbreak{}WITHOUT\_\allowbreak{}FIT\_\allowbreak{}\_\allowbreak{}NOT\_\allowbreak{}A\_\allowbreak{}CREDITATION\_\allowbreak{}\_\allowbreak{}GATE\_\allowbreak{}UNMOVED} --- o desfecho que a regra antiga daria, dito às claras.
\textbf{O achado do suporte e a emenda.} A tabela de erros publicada declara variância zero em 72 direções, que tocam 111 linhas (37 hospedeiras nas três âncoras); a colaboração resolve por pseudoinversa, e os pesos congelados em 06/11 são os dela e dão peso zero a essa componente. Em 11/11 o operador emendou a regra (D162): o ramo (1) deixou de recusar o vetor fora do suporte e passou a ser o portão 1 (a reprodução do número local publicado e a impressão digital). O literal novo só difere do antigo pela marca da emenda, e os ramos 2 a 6 são os mesmos (conferido neste rito). É emenda posterior à guarda, e o teste não é cego.
\textbf{O desenho ratificado (D159).} O primário é o modelo de erros publicado, com os 13 calibradores extras fora ($D$ de desenho $=6{,}71$); a linha das Cefeidas livres vai ao lado, sempre ($D=5{,}19$); o portão é de 90\,\% de poder de desenho no modelo publicado, com as linhas de estresse ao lado.
\textbf{A publicação antes da abertura.} O pré-registro foi publicado no commit \texttt{769fcb2697d8} (push 2026-11-12 10:45:40 UTC); a abertura começou 2026-11-12 10:46:34 UTC, uma única vez; o resultado foi publicado no commit \texttt{f24d954a3b90} (push 2026-11-12 10:55:51 UTC).
\textbf{O resultado (as duas linhas juntas).} No primário, sob o modelo de erros publicado, $Z$ contra o nulo $=6{,}11$ ($p$ unilateral $5{,}1\times10^{-10}$) e $Z$ contra $K=-0{,}61$, ramo (4); ao lado, no mesmo item, com as Cefeidas livres, $Z$ contra o nulo $=4{,}38$ e $Z$ contra $K=-0{,}81$, ramo (6), abaixo de 5. Em km/s/Mpc: o fundo $68{,}53$, a previsão da lei $74{,}27$, o lido $73{,}74$ e $73{,}35$. Retrodição sem ajuste; o teste não é cego; provado é validade, não creditação; o gate não se move.
\textbf{Os estresses (não mudam o veredito).} (i) Ruído por linha nas 111 linhas, com $\tau$ estimado do próprio vetor ($\hat\tau=0{,}0222$ mag, pós-observação): $Z$ contra o nulo $6{,}09$ e $4{,}37$. (ii) Âncora única: N4258 $3{,}89$, LMC $5{,}54$, MW $5{,}59$. (iii) Viés das Cefeidas variando por hospedeira, $B=0{,}02$ mag: $Z$ entre $5{,}56$ e $6{,}66$; $B=0{,}05$ mag: entre $4{,}73$ e $7{,}48$. Estas linhas são do primário sob o modelo de erros publicado, com as variações ditas; com as Cefeidas livres o dado lê $4{,}38$, abaixo de 5; retrodição sem ajuste; o teste não é cego; provado é validade, não creditação.
\textbf{O esperado pelos agregados, ao lado do observado ($Z$ contra o nulo).} primário (modelo publicado, 13 fora): esperado $5{,}83$, observado $6{,}11$; Cefeidas livres (13 fora): esperado $4{,}51$, observado $4{,}38$; âncora única N4258: esperado $4{,}06$, observado $3{,}89$; âncora única LMC: esperado $5{,}19$, observado $5{,}54$; âncora única MW: esperado $5{,}57$, observado $5{,}59$. É relato, não portão: o rito conferiu a conta. Estas linhas são do primário sob o modelo de erros publicado, com as variações ditas; com as Cefeidas livres o dado lê $4{,}38$, abaixo de 5; retrodição sem ajuste; o teste não é cego; provado é validade, não creditação.
\textbf{O que o teste não decide.} a creditação; a especificidade (o Z contra o nulo mede a tensão de Hubble; próprio da lei K é o tamanho, o Z contra K); se o orçamento de erros publicado está certo; a independência entre local e fundo; não é evidência independente das Fases 4, 5 e 9.
\textbf{As sete ressalvas.} (1) O ramo (1) foi emendado em 11/11/2026, depois de se saber que o vetor publicado o reprovava; pela função da v410 o desfecho era inconclusivo por sistemática. O teste não é cego. (2) O resultado do primário é previsível por dois agregados públicos, a menos de 0,18 desvio; o rito confere a conta, não descobre um número novo. É retrodição sem ajuste. (3) A linha das Cefeidas livres acompanha o resultado, no mesmo item. O esperado dela pelos agregados é 4,5, com 23 \% de chance de passar de 5. Nas três linhas de estresse de covariância o esperado também fica abaixo de 5. (4) O poder é de desenho no modelo publicado; as linhas de estresse vão ao lado. (5) A tabela de erros publicada declara zero em 72 direções; nelas, a maior componente do vetor publicado, na base do estimador congelado, é 0,0931 mag (o número depende da base; a norma, que não depende, sai no resultado). Isso não move o valor estimado, mas pesa no crédito do desvio. (6) O primário, o portão de 90 \% e as hipóteses foram escolhidos com a tabela de poder à vista; os agregados dos dois lados são públicos. (7) As tentativas anteriores e a sobreposição dos dados locais (seção 9).
\textbf{As tentativas anteriores.} A lei $K$ já foi confrontada com a razão de Hubble neste programa (Fase 4, Fase 5 e Fase 9 V2); os dados locais deste teste se sobrepõem aos daquelas, e ele não é evidência independente delas; as tentativas ficam contadas no Mapa de Rotas.
\textbf{E44} errata da gerência c1, em nome próprio (a seção “O teste pré-registrado da razão de Hubble com a rede de distâncias (H0DN $\times$ BG1)”, texto selado da v410 que o artigo da v411 traz igual): a seção descreve o estado da v410, de 09 a 11/11/2026. Onde ela diz “esta versão”, “nesta versão” e “hoje”, lê-se a v410. Onde diz que nenhum dado observado é lido, que o um.py não lê esse vetor e que a abertura é de versão futura, lê-se: assim era na v410; a versão futura é a v411, em que o vetor está embutido e o resultado é recalculado no rito. Onde trata o vetor fora do suporte como parede de montagem do dado, leia-se: é propriedade do modelo de erros publicado (o achado da gerência de 11/11 e a emenda, D162); não é erro de montagem. Onde condiciona a abertura a uma certificação prévia, às cegas, sob a covariância física, lê-se: a condição foi substituída em 11/11 pelo desenho ratificado (D159), com o primário no modelo de erros publicado e o portão de 90 \% de poder de desenho; a certificação não existiu e o teste aberto não é cego (ressalvas 1, 4 e 6). E onde textos de versões anteriores (o parágrafo “O que não se declara”, da v392, e a cadeia de máquina do registro estatuto\_da\_prova\_v391) dizem que $\beta$ na natureza a 5$\sigma$ aguarda instrumento, vale o ao lado da v411 no estatuto da prova. O texto selado não é reescrito; a errata vai ao lado.
